\chapter{El cierre dodecafásico de \texorpdfstring{\(\alpha\)}{alpha}}
\label{chap:alpha}
\index{alpha@\(\alpha\)}
\index{ciclo dodecafásico}

Esta sección establece el cierre dodecafásico mediante aritmética entera en
base mil. La construcción explicita las doce recurrencias, los cocientes
euclídeos propagados y las condiciones de frontera que determinan la salida.

\section{El selector terminal de los tres canales}
\label{sec:selector-terminal-alpha-rev6}

El cierre reúne tres generaciones ya construidas, pero no las identifica por
su nombre decimal. A profundidad terminal, los catálogos de colas compatibles
poseen cardinales
\[
 |\mathcal T_\pi|=196,\qquad
 |\mathcal T_e|=197,\qquad
 |\mathcal T_\varphi|=196.
\]
El objeto de selección es, por tanto,
\(\mathcal T_\pi\times\mathcal T_e\times\mathcal T_\varphi\), con
\[
 196\cdot197\cdot196=7\,567\,952
\]
ternas posibles.

La frontera de coordenadas fijada por el lector terminal es
\[
 q=(1,2,2,1,2,0)\in\mathbb F_3^6.
\]
La ley de canal \(\pi+e-\varphi\) posee vector de signos
\[
 \sigma=(1,1,-1)=(1,1,2)\in\mathbb F_3^3
\]
y recta dual no nula
\[
 \langle2\sigma\rangle
 =\langle(2,2,1)\rangle.
\]
El primer dato actúa sobre las columnas de la frontera; el segundo, sobre
sus filas. No son dos rasgos intercambiables.

\begin{theorem}[Selector matricial terminal de \(\alpha\)]
\label{thm:selector-terminal-alpha-rev6}
En el catálogo terminal anterior, las firmas locales de Gram, norma, peso y
distancia de Hamming reducen las \(7\,567\,952\) ternas a \(331\). Al exigir
el margen de columnas \(q\) quedan dos candidatas:
\[
 (120,197,157),\qquad(129,197,148).
\]
La condición de que la carga de filas pertenezca a la recta del canal
\(\langle\sigma\rangle\) selecciona exactamente la segunda. Sus bloques son
\[
 \boxed{
 b_\pi=021020,\qquad
 b_e=001220,\qquad
 b_\varphi=100210.}
\]
\end{theorem}

\begin{proof}
Las dos matrices terminales que sobreviven al margen de columnas son
\[
B_0=
\begin{pmatrix}
0&2&0&2&2&0\\
0&0&1&2&2&0\\
1&0&1&0&1&0
\end{pmatrix},
\qquad
B_1=
\begin{pmatrix}
0&2&1&0&2&0\\
0&0&1&2&2&0\\
1&0&0&2&1&0
\end{pmatrix}.
\]
Sus sumas de filas en \(\mathbb F_3\) son, respectivamente,
\[
 (0,2,0),\qquad(2,2,1).
\]
La primera no pertenece a
\(\langle(1,1,2)\rangle=\{0,(1,1,2),(2,2,1)\}\); la segunda sí, pues
\((2,2,1)=2\sigma\). Por tanto sólo \(B_1\) satisface simultáneamente la
frontera de coordenadas y la carga del canal. Sus tres filas son los bloques
mostrados y corresponden a la terna \((129,197,148)\).
\end{proof}

El teorema no obtiene \(q\) de la carga de filas: \(q\) es el dato
coordenado del lector terminal y \(\sigma\) es la orientación del canal.
Lo demostrado es que, una vez declaradas ambas fronteras del mismo objeto,
la rama real es única. Esta separación impide convertir una firma encontrada
a posteriori en una ley supuestamente derivada.

\begin{definition}[Dos lecturas del estado terminal enriquecido]
Sea
\[
 x_{\rm term}\in\mathcal X_{\rm TPK}^{\rm term,enr}
 \subseteq\mathcal X_{\rm TPK}^{[108]}.
\]
Este estado es la salida
\[
 x_{\rm term}
 =(\mathcal U_{108}\circ\cdots\circ\mathcal U_1)(x_{\rm seed})
\]
de la tubería completa de selección, transporte y actualización del TPK.
El selector anterior es la lectura
\[
 \pi_{\rm term}(x_{\rm term})
 =B_{\rm term}
 :=\begin{pmatrix}021020\\001220\\100210\end{pmatrix}.
\]
El registro orientado y el lector de los canales de \(90^\circ\) y
\(120^\circ\) producen, desde el libro del mismo estado,
\begin{equation}\label{eq:alpha-subespacio-firmado}
 \mathcal E_{108}^{90,120}\!
 \left(\mathcal R_{12}(x_{\rm term})\right)
 =(b^{90},b^{120},Q_{\rm TPK}),
\end{equation}
con
\[
\begin{aligned}
 b^{90}={}&(-495,-116,-684,108,601,-90,
             -173,-88,313,560,-397,461),\\
 b^{120}={}&(-425,-281,-481,671,-255,30,
              475,-543,680,251,6,-128),\\
 Q_{\rm TPK}={}&6263.
\end{aligned}
\]
Para \(r=1,2,3\), sean
\[
 B_r=\sum_{j=0}^{3}b^{120}_{r+3j},\qquad
 A_1=\frac{Q_{\rm TPK}+2B_1+B_2}{3},\qquad
 A_2=A_1-B_1,\qquad A_3=A_2-B_2,
\]
y \(d_{r,j}=b^{90}_{r+3j}\). El lector armónico \(\Pi_H\) forma los
bloques
\begin{equation}\label{eq:alpha-modulo-imagen-H12}
 \Pi_H^{(r)}
 =\bigl(A_r,d_{r,1}-d_{r,3},
              d_{r,0}+d_{r,2},d_{r,0}-d_{r,2}\bigr)
\end{equation}
y los intercala por columnas. La lectura firmada es, por tanto, la
composición producida
\begin{equation}\label{eq:alpha-Rsgn-definido}
 \boxed{
 R_{\rm sgn}:=
 \Pi_H\circ\mathcal E_{108}^{90,120}\circ\mathcal R_{12}:
 \mathcal X_{\rm TPK}^{\rm term,enr}
 \longrightarrow
 \mathcal U_{12}^{\rm sgn}:=\operatorname{im}R_{\rm sgn}.}
\end{equation}
Ni \(U\) ni \(K\) son coordenadas añadidas al dominio: \(U\) es la salida
de \eqref{eq:alpha-Rsgn-definido}, y \(K\) será su única preimagen integral
por Hadamard.
No se afirma una factorización
\(B_{\rm term}\to U_{12}^{\rm sgn}\): ambas salidas conservan componentes
distintas del mismo estado terminal.

Para \(N'\geq N\geq108\), los truncamientos actúan sólo sobre las historias
enriquecidas:
\begin{equation}\label{eq:alpha-truncamiento-firmado-preserva-KU}
 \tau_{N',N}:
 \mathcal X_{\rm TPK}^{[N'],\rm enr}
 \longrightarrow\mathcal X_{\rm TPK}^{[N],\rm enr}.
\end{equation}
Si \(R_{\rm sgn,N}\) denota la misma composición en profundidad \(N\), la
naturalidad es el cuadrado concreto
\begin{equation}\label{eq:alpha-cuadrado-naturalidad-Rsgn}
 \boxed{
 R_{\rm sgn,N}\circ\tau_{N',N}
 =R_{\rm sgn,N'}.}
\end{equation}
\end{definition}

\begin{theorem}[Construcción reversible de la coordenada firmada]
\label{thm:alpha-Rsgn-construido}
El compuesto \eqref{eq:alpha-Rsgn-definido}, aplicado al estado terminal
producido, da exactamente
\begin{equation}\label{eq:alpha-Rsgn-evaluacion}
 U=R_{\rm sgn}(x_{\rm term})
 =(2378,1406,2479,-452,998,-551,-668,-204,-371,-322,-28,-997).
\end{equation}
Su inversión por tres bloques de Hadamard de orden cuatro produce después
\[
 \boxed{
 K=(234,543,140,729,659,824,621,058,914,794,146,601).}
\]
La transformación es invertible; satisface el cuadrado
\eqref{eq:alpha-cuadrado-naturalidad-Rsgn} bajo los truncamientos
\eqref{eq:alpha-truncamiento-firmado-preserva-KU}; y es independiente de
\(\alpha\), CODATA y de toda cifra metrológica.
\end{theorem}

\begin{proof}
Las fórmulas \eqref{eq:alpha-subespacio-firmado} y
\eqref{eq:alpha-modulo-imagen-H12} dan
\[
 (B_1,B_2,B_3)=(972,-1073,101),\qquad
 (A_1,A_2,A_3)=(2378,1406,2479),
\]
y producen los tres bloques firmados
\[
\begin{aligned}
 U^{(1)}&=(2378,-452,-668,-322),\\
 U^{(2)}&=(1406,998,-204,-28),\\
 U^{(3)}&=(2479,-551,-371,-997).
\end{aligned}
\]
Su intercalado es \eqref{eq:alpha-Rsgn-evaluacion}. Como
\(H_4^2=4I_4\) y
\(\det(H_4\oplus H_4\oplus H_4)=(-16)^3=-4096\), se tiene
\begin{equation}\label{eq:alpha-Rsgn-inversa}
 K=\mathcal H_{12}^{-1}U
 =\frac14\mathcal H_{12}U
\end{equation}
después de deshacer la permutación de órbitas. En efecto,
\[
 H_4U^{(1)}=(936,2916,2484,3176),\quad
 H_4U^{(2)}=(2172,2636,232,584),
\]
y
\[
 H_4U^{(3)}=(560,3296,3656,2404).
\]
La división forzada por cuatro da las órbitas
\((234,729,621,794)\), \((543,659,058,146)\) y
\((140,824,914,601)\), cuyo intercalado es el sello declarado.

Como control independiente, los operadores cíclicos
\((D_3K)_m=K_m-K_{m+3}\) y
\((D_4K)_m=K_m-K_{m+4}\), junto con \(Q(K)=\sum_mK_m\), satisfacen
\begin{equation}\label{eq:alpha-extractor-90-120}
 \operatorname{rank}(D_3,D_4)=11,
 \qquad
 \operatorname{rank}(D_3,D_4,Q)=12.
\end{equation}
Por ello \((D_3K,D_4K,Q(K))\) reconstruye un único sello; además, el
cálculo recupera exactamente
\[
 D_3K=b^{90},\qquad D_4K=b^{120},\qquad Q(K)=Q_{\rm TPK},
\]
lo cual certifica el lector previo del registro, sin convertir sus salidas
en entradas. El truncamiento
\eqref{eq:alpha-truncamiento-firmado-preserva-KU} no reescribe las primeras
ciento ocho ventanas del libro; por ello ambos miembros de
\eqref{eq:alpha-cuadrado-naturalidad-Rsgn} leen el mismo registro y
producen \(U\). Ninguna de las matrices utilizadas contiene \(\alpha\) ni
un dato experimental.
\end{proof}

El calendario observable de 108 pasos es una proyección estrictamente más
pobre: reiniciar sus cursores cada 9 pasos borra precisamente las
coordenadas bidireccionales que el registro enriquecido conserva. Por
eso el registro crudo no sustituye al estado enriquecido y no se usa para
reconstruir \(R_{\rm sgn}\). El teorema de selección anterior construye
\(B_{\rm term}\); el teorema \ref{thm:alpha-Rsgn-construido} construye, de
manera separada, la segunda publicación del mismo estado.

\section{Estado dodecafásico y aplicaciones reversibles}

El teorema de esta sección recibe del
teorema~\ref{thm:alpha-Rsgn-construido} el vector firmado
\[
\Usgn=(2378,1406,2479,-452,998,-551,-668,-204,-371,-322,-28,-997).
\]
La ruta tipada demostrada a partir de esa entrada es
\[
\Usgn\longleftrightarrow K,
\]
\[
K\longleftrightarrow(D_3K,D_4K,Q),
\]
\[
(P,E,\Phi,K,c_{13}=0)
\longleftrightarrow A=\alpha_{12}.
\]
La evaluación directa no recibe \(\alpha\) ni un valor metrológico como
entrada. Recibe \(\Usgn\), las tres representaciones arquimedianas y la
condición de frontera. El
\cref{thm:selector-terminal-alpha-rev6} construye antes la selección de los
tres bloques terminales; la presente capa comienza con el vector firmado que
los transporta al cierre dodecafásico. La traza reducida de una sola
coordenada posicional no interviene en esta determinación.

\begin{remark}[Dos representaciones dodecafásicas de tipos distintos]
\[
U_{12}^{\rm cat}=U_6\Vert U_6
\]
es una palabra periódica del catálogo, mientras
\[
\Usgn=(2378,1406,2479,-452,998,-551,-668,-204,-371,-322,-28,-997)
\]
es un vector entero firmado del espacio completo de estados de TPK\@. El primero pertenece a una
representación de bloques y tiene periodo seis. El segundo conserva orientación, no es
una palabra base mil y no tiene ese periodo. No son intercambiables.
\end{remark}

\section{El vector canónico y las tres órbitas}

El vector dodecafásico es
\[
\boxed{
K=(234,543,140,729,659,824,621,058,914,794,146,601).}
\]
Su suma es
\[
Q(K)=6263.
\]
El paso tres divide \(\Cdoce\) en tres órbitas de cuatro sectores:
\[
\begin{aligned}
K^{(1)}&=(K_1,K_4,K_7,K_{10})=(234,729,621,794),\\
K^{(2)}&=(K_2,K_5,K_8,K_{11})=(543,659,058,146),\\
K^{(3)}&=(K_3,K_6,K_9,K_{12})=(140,824,914,601).
\end{aligned}
\]
Cada órbita recorre cuatro sectores separados por \(90^\circ\). Las etiquetas
cardinales norte--este--sur--oeste aparecen después de fijar origen y orientación; el orden
primario es \((m,m+3,m+6,m+9)\).

\section{Primera representación: transformada de Hadamard}

Sea
\[
H_4=
\begin{pmatrix}
1&1&1&1\\
1&1&-1&-1\\
1&-1&1&-1\\
1&-1&-1&1
\end{pmatrix}.
\]
La acción sobre las tres órbitas da
\[
\begin{aligned}
H_4K^{(1)}&=(2378,-452,-668,-322),\\
H_4K^{(2)}&=(1406,998,-204,-28),\\
H_4K^{(3)}&=(2479,-551,-371,-997).
\end{aligned}
\]
Al intercalar por columnas recuperamos
\[
\mathcal H_{12}(K)=\Usgn.
\]
Como
\[
H_4^2=4I_4,
\qquad
\det H_4=-16,
\]
la transformación global es invertible sobre \(\mathbb Q^{12}\):
\[
K^{(r)}=\frac14H_4U^{(r)},
\qquad
\det\mathcal H_{12}=(-16)^3=-4096\ne0.
\]
En la representación canónica, la inversa produce un único vector entero de
\([0,999]^{12}\).

\begin{theorem}[Estructura integral de la transformación de Hadamard]
\label{thm:smith-hadamard}
La forma normal de Smith del bloque \(H_4\), denotada por
\(\operatorname{SNF}(H_4)\), es
\[
 \operatorname{SNF}(H_4)=\operatorname{diag}(1,2,2,4).
\]
En consecuencia, para la suma directa de los tres bloques,
\[
 \operatorname{coker}\mathcal H_{12}
 \cong(\mathbb Z/2\mathbb Z)^6\oplus(\mathbb Z/4\mathbb Z)^3,
\]
y la imagen de \(\mathbb Z^{12}\) posee índice \(4096\). Un vector
\(u\in\mathbb Z^4\) pertenece a \(H_4\mathbb Z^4\) exactamente cuando
\[
 H_4u\equiv0\pmod 4.
\]
\end{theorem}

\begin{proof}
El máximo común divisor de las entradas de \(H_4\) es uno. Los menores de
orden dos son pares y alguno tiene valor absoluto dos; los menores de orden
tres son múltiplos de cuatro y alguno tiene valor absoluto cuatro; por
último, \(\lvert\det H_4\rvert=16\). Los divisores determinantes son,
por tanto, \(1,2,4,16\), de donde se obtienen los factores invariantes
\(1,2,2,4\). La suma directa de tres bloques repite estos factores y tiene
índice \(16^3=4096\). Puesto que \(H_4^{-1}=H_4/4\), la última
congruencia equivale a la integralidad de la preimagen.
\end{proof}

\section{Transportador ciclotómico relativo entre los sectores tres y cuatro}
\label{sec:transportador-ciclotomico-a34-rev8}

Sea \(\mathsf R\) la rotación regular del ciclo \(C_{12}\) sobre
\(\mathbb Q^{12}\), y sean
\[
 W_d^{\rm cyc}:=\ker\Phi_d(\mathsf R)
\]
sus sectores ciclotómicos. Los proyectores racionales sobre los sectores de
órdenes tres y cuatro pueden escribirse
\[
 \Pi_3=
 \frac14(I+\mathsf R^3+\mathsf R^6+\mathsf R^9)
 -\frac1{12}\sum_{k=0}^{11}\mathsf R^k,
\]
\[
 \Pi_4=
 \frac16(I-\mathsf R^2+\mathsf R^4-\mathsf R^6
              +\mathsf R^8-\mathsf R^{10}).
\]
Para fijar el orden de coordenadas, sea
\[
 \mathsf P_{\rm orb}(z_1,\ldots,z_{12})
 =\bigl((z_1,z_4,z_7,z_{10}),
        (z_2,z_5,z_8,z_{11}),
        (z_3,z_6,z_9,z_{12})\bigr).
\]
La transformada global usada arriba es, en el orden natural de \(C_{12}\),
\[
 \mathcal H_{12}
 :=\mathsf P_{\rm orb}^{-1}
   (H_4\oplus H_4\oplus H_4)\mathsf P_{\rm orb}.
\]
Es decir, \(H_4\oplus H_4\oplus H_4\) actúa en el orden de las tres
órbitas de paso tres, no sobre tres cuaternas contiguas del orden natural.
Definimos entonces el transportador relativo
\[
 A_{34}:=\left.\Pi_4\mathcal H_{12}\Pi_3
          \right|_{W_3^{\rm cyc}}:
 W_3^{\rm cyc}\longrightarrow W_4^{\rm cyc}.
\]

\begin{theorem}[Transportador ciclotómico dodecafásico]
\label{thm:transportador-ciclotomico-a34-rev8}
En las bases
\[
 u_1=(1,-1,0)^{\times4},\qquad
 u_2=(0,1,-1)^{\times4}
\]
de \(W_3^{\rm cyc}\), y
\[
 v_1=(1,0,-1,0)^{\times3},\qquad
 v_2=(0,1,0,-1)^{\times3}
\]
de \(W_4^{\rm cyc}\), se cumple
\[
 A_{34}u_1=\frac23(v_1-v_2),\qquad
 A_{34}u_2=\frac23(v_1+v_2),
\]
y, por tanto,
\[
 \boxed{\det A_{34}=\frac89}.
\]
En particular, \(A_{34}\) es un isomorfismo racional entre ambos sectores.
\end{theorem}

\begin{proof}
Se sustituyen las cuatro sucesiones periódicas anteriores en las expresiones
de \(\Pi_3,\Pi_4\) y en los tres bloques de Hadamard. La matriz de \(A_{34}\)
en las bases declaradas es
\[
 \begin{pmatrix}
  2/3&2/3\\
 -2/3&2/3
 \end{pmatrix},
\]
cuyo determinante es \(8/9\).
\end{proof}

La involución racional que intercambia \(W_3^{\rm cyc}\) y
\(W_4^{\rm cyc}\) mediante \(A_{34}\) y \(A_{34}^{-1}\), y fija
\(W_{12}^{\rm cyc}\), transporta la suma
\(W_4^{\rm cyc}\oplus W_{12}^{\rm cyc}\) a
\(W_3^{\rm cyc}\oplus W_{12}^{\rm cyc}\). En esta última suma, el polinomio
ciclotómico de la rotación es
\[
 \Phi_3(x)\Phi_{12}(x),
\]
que coincide con el polinomio del Coxeter de tipo \(E_6\). Ésta es una firma
espectral exacta y prepara el corredor geométrico posterior.

\begin{statusbox}[title={Transporte racional y continuación integral}]
El operador \(A_{34}\) realiza el transporte ciclotómico racional entre los
sectores de órdenes tres y cuatro. Al fijar \(W_{12}^{\rm cyc}\), sitúa la
dodecafase en el sector espectral exacto
\(\Phi_3\Phi_{12}\) de tipo \(E_6\). La continuación integral se construye
mediante dos fuentes posteriores y compatibles: el
\cref{thm:entrelazador-dodecafase-witt} proporciona la isometría
\(M_{12}\)-equivariante hacia el módulo de Witt, y el
\cref{thm:identificacion-isometrica-tres-proyectores} encadena los tres
proyectores de incidencia. Sobre esa realización integral, el
\cref{p12-83-c20-s6:thm:grafo-27-rectas-nuevo} construye el grafo de las
veintisiete rectas y el
\cref{p12-83-c20-s6:chap:holonomia-schlafli-e6-nuevo} identifica su acción de
Weyl. Así, \(A_{34}\) fija la etapa racional de la genealogía y los
entrelazadores de Witt--Schläfli realizan su prolongación reticular e
incidencial.
\end{statusbox}

La separación de signo es
\[
U_{12}^{+}
=(2378,1406,2479,0,998,0,0,0,0,0,0,0),
\]
\[
U_{12}^{-}
=(0,0,0,452,0,551,668,204,371,322,28,997),
\]
y
\[
\Usgn=U_{12}^{+}-U_{12}^{-}.
\]

\section{Segunda representación: diferencias transversales}

Con índices en \(\{1,\ldots,12\}\) y la convención cíclica
\[
 K_{m+r}:=K_{1+((m+r-1)\bmod12)},
\]
definimos
\[
(D_3K)_m=K_m-K_{m+3},
\qquad
(D_4K)_m=K_m-K_{m+4}.
\]
El extractor transversal dodecafásico es el operador
\[
 \mathcal E_{\mathrm{dod}}
 :=(D_3,D_4,Q).
\]
El subíndice ``dod'' declara su dominio. Los desplazamientos tres y cuatro
corresponden a separaciones angulares de \(90^\circ\) y \(120^\circ\) en el
ciclo orientado \(C_{12}\), pero \(\mathcal E_{\mathrm{dod}}\) no es el operador de
incidencia hexada--octada ni el núcleo angular utilizado posteriormente en
las torres.
Los datos son
\[
\begin{aligned}
D_3K={}&(-495,-116,-684,108,601,-90,\\
       &\qquad-173,-88,313,560,-397,461),\\[1mm]
D_4K={}&(-425,-281,-481,671,-255,30,\\
       &\qquad475,-543,680,251,6,-128),\\[1mm]
Q={}&6263.
\end{aligned}
\]
Si \(D_3v=D_4v=0\), el vector es constante en los pasos tres y cuatro. Como
\(\gcd(3,4)=1\), esos pasos conectan todo el ciclo y
\[
\ker D_3\cap\ker D_4=\langle\mathbf1\rangle.
\]
Por tanto
\[
\operatorname{rank}(D_3,D_4)=11.
\]
La carga total elimina la traslación constante, porque
\(Q(\mathbf1)=12\ne0\). Luego
\[
\boxed{\operatorname{rank}(D_3,D_4,Q)=12,}
\]
y
\[
K\longleftrightarrow(D_3K,D_4K,Q)
\]
es un sistema de coordenadas exacto.

\begin{theorem}[Forma normal integral del sistema transversal]
\label{thm:smith-dodecafase}
Sea
\[
 \mathcal D=(D_3,D_4,Q):
 \mathbb Z^{12}\longrightarrow\mathbb Z^{25}.
\]
Sus factores invariantes no nulos son
\[
 \operatorname{SNF}(\mathcal D)
 =\operatorname{diag}(\underbrace{1,\ldots,1}_{11},12),
\]
y, por consiguiente,
\[
 \operatorname{coker}\mathcal D
 \cong\mathbb Z^{13}\oplus\mathbb Z/12\mathbb Z.
\]
\end{theorem}

\begin{proof}
Las filas de \((D_3,D_4)\) son incidencias orientadas del grafo
\[
 \operatorname{Cay}
 \bigl(\mathbb Z/12\mathbb Z,\{\pm3,\pm4\}\bigr).
\]
Como los pasos tres y cuatro generan \(\mathbb Z/12\mathbb Z\), el grafo
es conexo.
Las incidencias de un árbol generador proporcionan once menores reducidos
unimodulares; los once primeros factores invariantes son, pues, iguales a
uno.

Todo menor de orden doce contiene la fila \(Q\), ya que el bloque de
incidencias tiene rango once. Al sumar las once primeras columnas a la
última, operación unimodular, la última columna se anula en las filas de
incidencia y vale doce en la fila \(Q\). Todos los menores máximos son
divisibles por doce. La elección de las once incidencias de un árbol da un
menor restante de valor absoluto uno y, por tanto, un determinante máximo
de valor absoluto doce. El último factor invariante es exactamente doce.
\end{proof}

La componente \(\mathbb Z/12\mathbb Z\) registra la condición de
compatibilidad integral de la carga total con el potencial reconstruido.
En particular, la unicidad racional probada por el rango se refina aquí a
una afirmación sobre la red integral y su única obstrucción torsional.

El observable firmado también puede reconstruirse desde estas coordenadas.
Para \(r=1,2,3\), sea
\[
B_r=\sum_{j=0}^{3}(D_4K)_{r+3j}.
\]
Se obtiene
\[
(B_1,B_2,B_3)=(972,-1073,101).
\]
Las sumas orbitales son
\[
A_1=\frac{Q+2B_1+B_2}{3}=2378,
\quad
A_2=A_1-B_1=1406,
\quad
A_3=A_2-B_2=2479.
\]
Si
\(d_{r,j}=(D_3K)_{r+3j}\) para \(j=0,1,2,3\), con la misma convención
cíclica, el bloque firmado de la órbita \(r\) es
\[
(A_r,d_{r,1}-d_{r,3},d_{r,0}+d_{r,2},d_{r,0}-d_{r,2}),
\]
exactamente el bloque Hadamard correspondiente.

\begin{remark}[Codominios de \(\mathcal H_{12}\) y
\(\mathcal E_{\mathrm{dod}}\)]
La escritura correcta es
\[
\mathcal H_{12}(K)=\Usgn,
\qquad
\mathcal E_{\mathrm{dod}}(K)=(D_3K,D_4K,Q).
\]
Ambas representaciones son compatibles porque reconstruyen el mismo vector,
pero no se ha definido aquí un diagrama entre sus codominios. Sus fórmulas y
codominios son distintos. Compartir los pasos \(3/4\) o los índices \(90/120\) no
autoriza a identificarlas.
\end{remark}

\section{Las tres coordenadas arquimedianas}

Los vectores de doce tríadas son
\[
\begin{aligned}
P={}&(141,592,653,589,793,238,462,643,383,279,502,884),\\
E={}&(718,281,828,459,045,235,360,287,471,352,662,497),\\
\Phi={}&(618,033,988,749,894,848,204,586,834,365,638,117).
\end{aligned}
\]
En forma concatenada:
\[
\begin{aligned}
\iota(P)&=141592653589793238462643383279502884,\\
\iota(E)&=718281828459045235360287471352662497,\\
\iota(\Phi)&=618033988749894848204586834365638117.
\end{aligned}
\]
\(P\) es la coordenada decimal compartida por las cinco regiones de
\(\pi\); el cálculo no afirma que esas regiones sean idénticas.

\section{Recurrencia de acarreo en base mil}

Fijamos
\[
c_{13}=0
\]
y calculamos de derecha a izquierda, \(m=12,11,\ldots,1\):
\[
s_m=P_m+E_m-\Phi_m-K_m+c_{m+1},
\]
\[
A_m=s_m\bmod1000,
\qquad
c_m=\frac{s_m-A_m}{1000},
\qquad0\le A_m<1000.
\]
La tabla se imprime en orden creciente de índice, pero cada acarreo procede de la fila
inferior.

\begingroup
\scriptsize
\setlength{\tabcolsep}{3.3pt}
\begin{longtable}{@{}r|rrrr|rr|rr@{}}
\caption{Cierre dodecafásico completo en base mil.}\label{tab:acarreo-alpha}\\
\toprule
\(m\)&\(P_m\)&\(E_m\)&\(\Phi_m\)&\(K_m\)&\(c_{m+1}\)&\(s_m\)&\(A_m\)&\(c_m\)\\
\midrule
\endfirsthead
\toprule
\(m\)&\(P_m\)&\(E_m\)&\(\Phi_m\)&\(K_m\)&\(c_{m+1}\)&\(s_m\)&\(A_m\)&\(c_m\)\\
\midrule
\endhead
1&141&718&618&234&0&7&007&0\\
2&592&281&033&543&0&297&297&0\\
3&653&828&988&140&\(-1\)&352&352&0\\
4&589&459&749&729&\(-1\)&\(-431\)&569&\(-1\)\\
5&793&045&894&659&\(-2\)&\(-717\)&283&\(-1\)\\
6&238&235&848&824&\(-1\)&\(-1200\)&800&\(-2\)\\
7&462&360&204&621&0&\(-3\)&997&\(-1\)\\
8&643&287&586&058&\(-1\)&285&285&0\\
9&383&471&834&914&\(-1\)&\(-895\)&105&\(-1\)\\
10&279&352&365&794&0&\(-528\)&472&\(-1\)\\
11&502&662&638&146&0&380&380&0\\
12&884&497&117&601&0&663&663&0\\
\bottomrule
\end{longtable}
\endgroup

Los acarreos son
\[
(c_1,\ldots,c_{13})
=(0,0,0,-1,-1,-2,-1,0,-1,-1,0,0,0).
\]
Los valores negativos de \(s_m\) son préstamos ordinarios en base mil. La
división euclídea fija resto y cociente de forma única.

La salida es
\[
\boxed{
A=(007,297,352,569,283,800,997,285,105,472,380,663),}
\]
y por tanto
\[
\boxed{
\begin{aligned}
\alpha_{12}
&:=\frac{7297352569283800997285105472380663}{10^{36}},\\
&=0.007297352569283800997285105472380663.
\end{aligned}}
\]
Este racional es la ventana dodecafásica inicial de la sección compatible
generada por el transporte TPK. Si
\(\rho_{N',N}:C_{N'}^\alpha\to C_N^\alpha\) son sus truncamientos, la
salida completa es
\[
 \boxed{
 \alpha_{\mathrm{HMT}}
 :=\varprojlim_N\alpha_N,\qquad
 \rho_{N',N}(\alpha_{N'})=\alpha_N,\qquad
 \alpha_{36}=\alpha_{12}.}
\]
El retorno de fase prolonga la memoria y estrecha los cilindros
supervivientes a profundidad arbitraria; no reinicia el estado ni termina la
construcción en treinta y seis cifras.

\section{La identidad afín completa}

Coordenada a coordenada,
\[
\boxed{
P_m+E_m-\Phi_m-K_m+c_{m+1}
=A_m+1000c_m.}
\]
Si
\[
C=(c_1,\ldots,c_{12}),
\qquad
C^+=(c_2,\ldots,c_{13}),
\]
la forma vectorial correcta es
\[
\boxed{P+E-\Phi-K+C^+-1000C=A.}
\]
Para evitar identificar las colas finitas con las constantes completas,
definimos
\[
 \tau_{36}(x)=10^{-36}\left\lfloor10^{36}
 \bigl(x-\lfloor x\rfloor\bigr)\right\rfloor,
 \qquad
 \kappa(K)=10^{-36}\iota(K).
\]
La notación decimal exacta de la identidad concatenada es entonces
\[
 \tau_{36}(\pi)+\tau_{36}(e)-\tau_{36}(\varphi)-\kappa(K)
 =\alpha_{12}.
\]
Sin el truncamiento explícito, la inmersión \(\kappa\) y el cociclo
\(C^+-1000C\), la igualdad coordenada sería falsa.

Definimos
\[
\iota(v_1,\ldots,v_{12})
=\sum_{m=1}^{12}v_m1000^{12-m}.
\]
Al multiplicar cada recurrencia por \(1000^{12-m}\) y sumar, los acarreos
telescopan:
\[
\iota(P)+\iota(E)-\iota(\Phi)-\iota(K)-\iota(A)
=1000^{12}c_1-c_{13}.
\]
Como \(c_1=c_{13}=0\):
\[
\boxed{
\iota(P)+\iota(E)-\iota(\Phi)-\iota(K)=\iota(A).}
\]
Los dos enteros restantes son
\[
\iota(K)=234543140729659824621058914794146601,
\]
\[
\iota(A)=007297352569283800997285105472380663.
\]

\section{Inversión y equivalencia de las tres representaciones}

La recurrencia se invierte por
\[
K_m=P_m+E_m-\Phi_m+c_{m+1}-A_m-1000c_m,
\]
y la representación entera concatenada por
\[
\iota(K)=\iota(P)+\iota(E)-\iota(\Phi)-\iota(A).
\]
El diagrama completo es
\[
\boxed{
K\longleftrightarrow\Usgn,
\qquad
K\longleftrightarrow(D_3K,D_4K,Q),
\qquad
K\longleftrightarrow A\mid(P,E,\Phi,c_{13}=0).}
\]
Las dos primeras son representaciones lineales; la tercera, una aplicación
afín reversible en la fibra
condicionada por \(P,E,\Phi\) y la frontera.

\begin{remark}[Conmutatividad del cierre]
La aplicación directa obtiene \(K\) desde la representación firmada y produce
\(A\). La fórmula inversa demuestra que la salida y la frontera recuperan el
mismo vector sin libertad residual. La invertibilidad no constituye una
evidencia estadística independiente; la conmutatividad de las tres
representaciones es una propiedad algebraica del cierre.
\end{remark}

La unicidad se expresa en tres niveles de la misma construcción:
\begin{enumerate}
  \item fijados el vector firmado, \(P,E,\Phi\) y la frontera, el vector y la
  salida son únicos;
  \item las representaciones son reversibles y conmutativas;
  \item la genealogía completa
  \(\mathrm{APP}\to\mathrm{TRIT}\to\mathrm{TPK}\) produce esas primitivas
  antes del cierre local, y sus truncamientos compatibles prolongan la
  identidad a la sección coinductiva \(\alpha_{\mathrm{HMT}}\).
\end{enumerate}

\section{Condición de frontera dodecafásica}

La primera ventana dodecafásica usa \(c_{13}=0\) y termina en
\[
\ldots|380|663.
\]
Esta frontera forma parte del dominio de la aplicación afín reversible. La
prolongación no introduce otro sistema: itera el mismo transporte graduado
del TPK sobre
\[
 w_6\to w_{12}\to w_{18}\to w_{24}\to w_{30}\to R_{36}\to G_9
\]
y construye las ventanas sucesivas con retorno de fase, avance de memoria y
acarreo terminal compatible. El
\cref{sec:l9s102:alpha-dos-vias} demuestra que los cuadrados de truncamiento
de esta vía y del lector analítico conmutan y que sus cilindros
supervivientes poseen intersección unitaria.

\subsection{Clase integral del acarreo}

La condición terminal puede compararse con el problema periódico. Sea
\(S\) el desplazamiento cíclico sobre \(\mathbb Z^{12}\) y, para una base
entera \(b\geq2\), definamos
\[
 \partial_b=S-bI.
\]
La identidad afín periódica adopta la forma
\[
 A=P+E-\Phi-K+\partial_bC.
\]
Para todo \(h\in\mathbb Z^{12}\), la transformación
\[
 K\longmapsto K+\partial_bh,
 \qquad
 C\longmapsto C+h
\]
conserva \(A\). Así, la elección de acarreos es una elección de
representante dentro de una clase integral.

\begin{proposition}[Cociente periódico del acarreo]
\label{prop:cociente-periodico-acarreo}
En un ciclo de longitud doce,
\[
 \operatorname{coker}(S-bI)\cong
 \mathbb Z/(b^{12}-1)\mathbb Z.
\]
La concatenación en base \(b\), tomada módulo \(b^{12}-1\), determina la
clase periódica.
\end{proposition}

\begin{proof}
El módulo cíclico se identifica con
\(\mathbb Z[t]/(t^{12}-1)\). Imponer la relación \(t=b\) produce
\(\mathbb Z/(b^{12}-1)\mathbb Z\). La evaluación del polinomio de
coeficientes es la concatenación en base \(b\), con la orientación fijada
por el orden de los doce sectores.
\end{proof}

El cierre de \(\alpha\) no emplea el cociente periódico: utiliza una cadena
orientada con extremo derecho y exige \(c_{13}=0\). La división euclídea
selecciona entonces un único representante. Esta diferencia explica por
qué la frontera abierta forma parte de la definición de la aplicación y no
puede reemplazarse por una identificación cíclica de los acarreos.

\section{Comparación con la referencia metrológica}

Después de la generación, la aritmética entera de la primera ventana
verifica el reconocimiento
\[
\boxed{
\alpha_{12}
=10^{-36}
\left\lfloor\frac{10^{36}}{137.035999084}\right\rfloor.}
\]
Su inversa es
\[
\alpha_{12}^{-1}
=137.03599908400000000000000000000001066588\ldots,
\]
y reproduce exactamente la cadena central publicada en el ajuste de 2018 del
Committee on Data for Science and Technology (CODATA) \cite{Tiesinga2021CODATA}:
\[
\boxed{137.035999084.}
\]
No es una coincidencia de coma flotante. Las treinta y seis cifras de
\(\alpha_{12}\) constituyen el truncamiento exacto del recíproco de la
cadena central impresa. Esta comprobación finita no define
\(\alpha_{\mathrm{HMT}}\) ni limita su profundidad: la sección coinductiva
ya producida determina todas sus ventanas compatibles, mientras CODATA sólo
reconoce después el prefijo que su resolución experimental permite
contrastar.

El reconocimiento metrológico no forma parte de la recurrencia. El resultado
matemático es la identidad exacta de las dos vías
\[
 \alpha_A=\alpha_B=\alpha_{\mathrm{HMT}},
\]
cuyas ventanas dodecafásicas son determinadas por el vector firmado, las
tres palabras generadas y las fronteras compatibles. La tabla de CODATA
sólo compara después una salida ya fijada.

La comparación con ajustes metrológicos posteriores es un control externo
separado. La salida HMT queda congelada y, por ello, puede ser contrastada sin
modificar el generador.

\section{Independencia de las coordenadas y rigidez}

Las pruebas del cierre distinguen modificaciones profundas de perturbaciones
locales:
\begin{center}
\small
\begin{tabularx}{\textwidth}{@{}lYY@{}}
\toprule
Variante & Resultado & Qué controla \\
\midrule
protocolo canónico derecha--izquierda & igualdad exacta & referencia \\
usar \(U_6\Vert U_6\) como \(K\) & falla macroscópicamente & distinción de tipos \\
rotar \(K\) & rompe el cierre & origen dodecafásico \\
reflejar \(K\) & invierte el borde & orientación \\
acarreo izquierda--derecha & cambia la palabra & sentido del transporte \\
perturbar \(K_i\) en una unidad & falla la palabra exacta & rigidez profunda \\
\bottomrule
\end{tabularx}
\end{center}
Una perturbación local puede alterar sólo una cifra profunda; sigue siendo una
refutación de la igualdad exacta. Las variantes estructurales producen fallos
mayores y muestran que el cierre no es invariante bajo órdenes arbitrarios.

\begin{remark}[Estructura del cierre dodecafásico]
El cierre de \(\alpha\) comprende un ciclo orientado, dos representaciones
exactas del vector \(K\), doce entradas tipadas, una recurrencia de derecha a
izquierda, trece acarreos, una condición de frontera y una identidad entera
telescópica. Estos datos determinan la salida de manera reproducible.
\end{remark}
